%!TeX root=csp.tex
\section{Standing conventions}\label{sec:conventions}

Zhuk's papers carry their hypotheses in prose: ``in this paper we assume that
every algebra is a finite idempotent algebra having a WNU term operation''.
A formalization cannot carry a hypothesis in prose, and an informal reader
cannot audit which results actually consume which hypothesis. We therefore fix
the hypotheses as a labelled convention and record, for each one, where it is
consumed.

\begin{convention}[Standing hypotheses]\label{conv:standing}
Throughout, $n\ge 3$ is a fixed integer and:
\begin{enumerate}\alphlabels\tightlist
\item\label{conv:finite} every algebra is \emph{finite} and has a
  \emph{nonempty} universe;
\item\label{conv:onebasic} every algebra has exactly one basic operation, of
  arity $n$, written $w$;
\item\label{conv:idem} $w$ is \emph{idempotent}: $w(x,\dots,x)=x$;
\item\label{conv:wnu} $w$ is a \emph{weak near-unanimity} operation:
  $w(y,x,\dots,x)=w(x,y,x,\dots,x)=\dots=w(x,\dots,x,y)$ for all $x,y$;
\item\label{conv:special} $w$ is \emph{special}:
  $w(x,\dots,x,y)=w\bigl(x,\dots,x,w(x,\dots,x,y)\bigr)$ for all $x,y$.
\end{enumerate}
The class of algebras satisfying \ref{conv:finite}--\ref{conv:special} is
written $\Vn$. Since only finite algebras are considered, $\Vn$ is not a
variety.
\end{convention}

\begin{remark}[Why one basic operation]\label{rem:one-basic}
Restricting to a single basic operation is not a loss. The dichotomy hypothesis
is that \emph{some} weak near-unanimity operation preserves $\Gamma$; one may
then discard the original signature and work in the algebra
$(A;w)$ whose only basic operation is that one, because every relation of
$\Gamma$ is a subuniverse of a power of $(A;w)$ and that is all the argument
uses. Passing from a WNU $w$ to a \emph{special} one is
Lemma~\ref{lem:special-wnu}, and costs an increase of arity from $n$ to
$n^{n!}$.

The gain is large for a formalization: a term is a finite tree with $n$-ary
branching over a single symbol, so ``term operation'' needs no signature
bookkeeping, and $\Clo(\mathbf A)$ is the clone generated by one operation.
\end{remark}

\begin{hazard}[Where each hypothesis is consumed]\label{haz:hyp-audit}
\ref{conv:finite} is used in every induction on $|A|$ and in every
minimality argument (a minimal element of a nonempty family of subsets
exists); it cannot be dropped anywhere. \ref{conv:idem} is what makes
singletons subuniverses and makes $\Sg$ of a nonempty set nonempty.
\ref{conv:wnu} is consumed only through the results it implies --- the
existence of absorption-theoretic structure --- and never directly in
\S\S\ref{sec:instances}--\ref{sec:main-statements}. \ref{conv:special} is used
where a unary polynomial must be idempotent under iteration, chiefly in the
theory of $\Sg$-closures of two-element sets. A formalization should carry
\ref{conv:finite}--\ref{conv:special} as typeclass instances on a structure and
not as ambient hypotheses, so that the audit is mechanical.
\end{hazard}

\begin{convention}[The empty set]\label{conv:empty}
A \emph{subuniverse} of $\mathbf A$ is a subset closed under $w$; the empty set
is a subuniverse under this definition, and we allow it. A
\emph{subalgebra} is a nonempty subuniverse; $\mathbf B\le\mathbf A$ always
carries $B\neq\varnothing$.

The relations $\subt{T}$, $\subte{T}$ and $\lll$ of \S\ref{sec:types} are
defined only between \emph{nonempty} sets; in particular
$\varnothing\lll A$ never holds. Where a construction can produce the empty
set --- a projection of an intersection, most often --- the dotted variants
$\dsubt{T}$, $\dsubte{T}$, $\dlll$ are used, and they mean ``the stated relation
holds, or the left-hand side is empty''.
\end{convention}

\begin{hazard}[The dot is not decoration]\label{haz:dot}
Conflating $\lll$ with $\dlll$ is the single most common way to misread
\S\ref{sec:types}. Every propagation statement of \S\ref{sec:properties} that
projects or intersects produces a dotted conclusion, and every statement that
consumes a $\lll$ hypothesis needs the undotted one. In a formalization the two
should be distinct constants, with the dotted one \emph{defined} as
\texttt{C = }$\varnothing$\texttt{ }$\vee$\texttt{ C }$\lll$\texttt{ B}, so that the case split is forced
at every use site rather than being resolved by the reader's judgement.
\end{hazard}

\begin{convention}[Relations, and indices]\label{conv:relations}
A relation is a subset of a product $A_{1}\times\dots\times A_{m}$ of universes,
indexed by $[m]=\{1,\dots,m\}$; $\mathbf R\le\mathbf A_{1}\times\dots\times
\mathbf A_{m}$ means additionally that $R$ is a subuniverse of the product
algebra. We write $\proj_{i_{1},\dots,i_{s}}(R)$ for the projection onto the
listed coordinates and $\mathbf R\sdle\mathbf A_{1}\times\dots\times\mathbf
A_{m}$ when every $\proj_{i}(R)=A_{i}$.

Products are indexed by \emph{arbitrary finite index sets}, not only by initial
segments of $\mathbb N$. Several constructions in \S\ref{sec:main-statements}
form a product over the variable set of an instance, or over a set of
subuniverses, and reindexing those to $[m]$ by hand is exactly the kind of
bookkeeping that a formalization pays for repeatedly and that buys nothing.
\end{convention}
